\section{Algorithms and the numerical claims}\label{sec:computing}

\begin{theorem}[All values and an isolated value]\label{thm:algorithm}
All the values through $N$ take $O(N)$ word operations and $O(N)$ words. For one value at $N\ge1$, the quotient recurrence takes $O(N^{2/3})$ word operations and $O(N^{2/3})$ words with the cutoff below.
\end{theorem}
\begin{proof}
For the table maintain a list of the primes found and a Boolean composite array initialized to false. Set $\mu(1)=1$. In increasing order of $i\ge2$, if $i$ is unmarked append it to the prime list and set $\mu(i)=-1$. For each prime $p$ in list order form $j=ip$. Stop if $j>N$. Mark $j$ composite. If $p\mid i$, set $\mu(j)=0$ and end this inner loop. Otherwise set $\mu(j)=-\mu(i)$ and continue.

Every composite $j$ has the unique generating pair $(i,p)$ with $p$ its least prime factor and $i=j/p$. Its least factor satisfies $p\le i$, and no prime below $p$ divides $i$. The pass at $i$ therefore reaches $p$. Conversely each product generated has $p$ as least factor, by the stopping condition. Thus the prime test on an unmarked $i$ is justified inductively, each composite is marked just once, and the update for $\mu$ follows from its prime exponents. There is at most one extra terminating test per outer iteration. A final prefix sum gives $M$.

For the isolated-value method let $b=\lceil\sqrt[3]N\rceil$, $y=\min(N,b^2)$ and $K=\lfloor N/(y+1)\rfloor\le b$. First sieve to $y$. For $i=K,K-1,\ldots,1$ compute the value at $v=\lfloor N/i\rfloor$. In its recurrence, starting at $a=2$, one whole quotient block has
\[
 u=\lfloor v/a\rfloor,\qquad e=\lfloor v/u\rfloor,\qquad
 \text{contribution }(e-a+1)M(u).
\]
Advance to $e+1$. For $u\le y$ use the sieve. For $u>y$, use the previously computed entry at index $ia$, since $u=\lfloor N/(ia)\rfloor$ and $i<ia\le K$. This proves the reverse-loop invariant and correctness, including the case $K=0$ when the sieve already gives the answer.

There are at most $2\lfloor\sqrt{N/i}\rfloor$ blocks at $i$. To sum this bound count the pairs $(i,j)$ with $1\le i\le K$, $j\ge1$ and $ij^2\le N$. Those with $j\le b$ contribute at most $b^2$ and the others at most $\sum_{j>b}N/(j(j-1))=N/b\le b^2$. This proves the time bound with integer cutoffs throughout. A running total in a recurrence has magnitude at most $1+N^2$. Indices and products fit in a constant multiple of $\log(2N)$ bits. The operation model allows indexed word reads and writes and integer division. With schoolbook word arithmetic in a random-access bit-cost model one can charge $O(\log^2(2N))$ per operation, giving the $O_\varepsilon(N^{2/3+\varepsilon})$ exponent. The operation-count claim concerns this algorithm. Other methods achieve smaller space bounds \cite{ht}.
\end{proof}

\paragraph{An independent formula for verification.}
Let $r=\lfloor\sqrt N\rfloor$. Using $\mu*\mu*\mathbf1=\mu$, sum the triple convolution over $abc\le N$ with weight $\mu(a)\mu(b)$. The regions $a>r$ and $b>r$ are disjoint. For fixed $a>r$ the sum on $(b,c)$ is $\sum_{m\le N/a}\sum_{b\mid m}\mu(b)=1$. Thus these two regions together contribute $2(M(N)-M(r))$. Summing $c$ in the remaining region and rearranging gives the identity used in the second computation \cite{ht}:
\[
 M(N)=2M(r)-\sum_{a,b\le r}\mu(a)\mu(b)\left\lfloor\frac{N}{ab}\right\rfloor.
\]

\paragraph{Exact finite evaluations.}
Both isolated-value algorithms yield the following ten integers.
\begin{center}
\begin{tabular}{@{}lrrrrrrrrrr@{}}
\toprule
$j$&0&1&2&3&4&5&6&7&8&9\\
$M(10^j)$&1&$-1$&1&2&$-23$&$-48$&212&1037&1928&$-222$\\
\bottomrule
\end{tabular}
\end{center}
For $N=1,\ldots,12$ the values are
$1,0,-1,-1,-2,-1,-2,-2,-2,-1,-2,-2$. The first ten positive zero arguments are
$2,39,40,58,65,93,101,145,149,150$. The full table to $428$ recovers the displayed initial $143$ terms, the displayed zero sequence, and the seven level sets in the two articles over their displayed ranges. The per-integer table is in \texttt{verification/initial-values.csv}. The program lists all the level hits over the required interval.

More sharply, on the integers $2\le n\le10\,000$ the maximum of $M(n)^2/n$ is $4/5$, first attained at $5$. The program verifies this by comparing integer cross products using two independent definitions of $\mu$. It uses the sieve and independently performs the update $\mu(m)\mathrel{-}=\mu(d)$ on each proper multiple $m$ of successive $d$, starting from the array $\delta$. Before visiting $d$, all contributions of its proper divisors have already been subtracted. This proves the update by the divisor identity. On each cell $[n,n+1)$ the largest ratio $M(n)^2/x$ occurs at $n$. On $(1,2)$ it is $1/x<1$. Together these give Mertens's historical strict inequality on $1<x\le10\,000$. For each finite calculation the proof rule in iPRA is evaluation of closed PR terms by their defining equations. The bundled scripts are reproducible evaluations, with the mathematical justifications above.

\paragraph{The disproof as a logical statement.}
The 1985 result of Odlyzko and te Riele gives excursions on both sides of the square-root bounds \cite{ot}. Its strict-excess conclusion implies $\exists n>1\ [M(n)^2>n]$ by taking an integer part. At the metamathematical level every true closed existential statement with a PR witness test has an iPRA proof obtained by evaluation at a witnessing numeral followed by existential introduction. This observation identifies the proof-theoretic status of the individual existential statement. It supplies neither that numeral nor a uniform oscillation theorem. The large oscillations in the 1985 paper and the effective bound of Pintz quoted in \cite{ht} remain external inputs here.
